% Options for packages loaded elsewhere
\PassOptionsToPackage{unicode}{hyperref}
\PassOptionsToPackage{hyphens}{url}
\PassOptionsToPackage{dvipsnames,svgnames,x11names}{xcolor}
\documentclass[
  british,
  11pt,
  a4paper,
]{article}
\usepackage{xcolor}
\usepackage[margin=23mm]{geometry}
\usepackage{amsmath,amssymb}
\setcounter{secnumdepth}{-\maxdimen} % remove section numbering
\usepackage{iftex}
\ifPDFTeX
  \usepackage[T1]{fontenc}
  \usepackage[utf8]{inputenc}
  \usepackage{textcomp} % provide euro and other symbols
\else % if luatex or xetex
  \usepackage{unicode-math} % this also loads fontspec
  \defaultfontfeatures{Scale=MatchLowercase}
  \defaultfontfeatures[\rmfamily]{Ligatures=TeX,Scale=1}
\fi
\usepackage{lmodern}
\ifPDFTeX\else
  % xetex/luatex font selection
\fi
% Use upquote if available, for straight quotes in verbatim environments
\IfFileExists{upquote.sty}{\usepackage{upquote}}{}
\IfFileExists{microtype.sty}{% use microtype if available
  \usepackage[]{microtype}
  \UseMicrotypeSet[protrusion]{basicmath} % disable protrusion for tt fonts
}{}
\makeatletter
\@ifundefined{KOMAClassName}{% if non-KOMA class
  \IfFileExists{parskip.sty}{%
    \usepackage{parskip}
  }{% else
    \setlength{\parindent}{0pt}
    \setlength{\parskip}{6pt plus 2pt minus 1pt}}
}{% if KOMA class
  \KOMAoptions{parskip=half}}
\makeatother
\usepackage{longtable,booktabs,array}
\newcounter{none} % for unnumbered tables
\usepackage{calc} % for calculating minipage widths
% Correct order of tables after \paragraph or \subparagraph
\usepackage{etoolbox}
\makeatletter
\patchcmd\longtable{\par}{\if@noskipsec\mbox{}\fi\par}{}{}
\makeatother
% Allow footnotes in longtable head/foot
\IfFileExists{footnotehyper.sty}{\usepackage{footnotehyper}}{\usepackage{footnote}}
\makesavenoteenv{longtable}
\usepackage{graphicx}
\makeatletter
\newsavebox\pandoc@box
\newcommand*\pandocbounded[1]{% scales image to fit in text height/width
  \sbox\pandoc@box{#1}%
  \Gscale@div\@tempa{\textheight}{\dimexpr\ht\pandoc@box+\dp\pandoc@box\relax}%
  \Gscale@div\@tempb{\linewidth}{\wd\pandoc@box}%
  \ifdim\@tempb\p@<\@tempa\p@\let\@tempa\@tempb\fi% select the smaller of both
  \ifdim\@tempa\p@<\p@\scalebox{\@tempa}{\usebox\pandoc@box}%
  \else\usebox{\pandoc@box}%
  \fi%
}
% Set default figure placement to htbp
\def\fps@figure{htbp}
\makeatother
\ifLuaTeX
\usepackage[bidi=basic,shorthands=off]{babel}
\else
\usepackage[bidi=default,shorthands=off]{babel}
\fi
\ifLuaTeX
  \usepackage{selnolig} % disable illegal ligatures
\fi
\setlength{\emergencystretch}{3em} % prevent overfull lines
\providecommand{\tightlist}{%
  \setlength{\itemsep}{0pt}\setlength{\parskip}{0pt}}
\usepackage{microtype}
\usepackage{fancyhdr}
\usepackage{titlesec}
\usepackage{xurl}
\usepackage{float}
\usepackage{caption}
\floatplacement{figure}{H}
\pagestyle{fancy}
\fancyhf{}
\fancyhead[L]{\small J-Lens in Ouro}
\fancyhead[R]{\small Working draft v2}
\fancyfoot[C]{\small\thepage}
\renewcommand{\headrulewidth}{0.3pt}
\setlength{\headheight}{14pt}
\setlength{\parindent}{0pt}
\titleformat{\section}{\large\bfseries}{}{0pt}{}
\titleformat{\subsection}{\normalsize\bfseries}{}{0pt}{}
\titlespacing*{\section}{0pt}{14pt plus 2pt}{6pt}
\titlespacing*{\subsection}{0pt}{10pt plus 2pt}{4pt}
\captionsetup{font=small,labelfont=bf}
\setkeys{Gin}{width=\linewidth}
\usepackage{bookmark}
\IfFileExists{xurl.sty}{\usepackage{xurl}}{} % add URL line breaks if available
\urlstyle{same}
% fallback for those not using the hyperref driver hyperxmp:
\makeatletter
\@ifundefined{xmpquote}{\newcommand{\xmpquote}[1]{#1}}{}
\makeatother
\hypersetup{
  pdftitle={J-Lens in Ouro: Early-Pass Deficits and a Confirmed Late-Pass Advantage},
  pdfauthor={Jan Kirin},
  pdflang={en-GB},
  colorlinks=true,
  linkcolor={Maroon},
  filecolor={Maroon},
  citecolor={Blue},
  urlcolor={Blue},
  pdfcreator={LaTeX via pandoc}}

\title{J-Lens in Ouro: Early-Pass Deficits and a Confirmed Late-Pass
Advantage}
\author{Jan Kirin}
\date{Working draft v2 · 11 September 2026}

\begin{document}
\maketitle

\subsection{Abstract}\label{abstract}

Does a Jacobian-based readout reveal intermediate content that a model's
own output head misses during recurrent computation? I compare J-Lens
with the raw logit lens in frozen Ouro-2.6B, which applies 48 shared
blocks over four recurrent passes. On multihop questions the raw lens
was stronger in passes 1--3, while layerwise analysis found a positive
J-Lens region in pass 4 whose location and size five independent
calibration fits reproduced on the discovery set. I then froze one
estimator (fit01) and the band, physical layers 26--37 of pass 4, before
evaluating 160 new two-hop questions with 80 new intermediate concepts.
In that band J-Lens improves excess hit@10 by \textbf{23.19 percentage
points} (95\% dependency-group bootstrap interval
\textbf{16.29--32.90}); intended-concept recovery is 37.92\% against
14.43\% for the raw lens. The same physical band shows the opposite sign
in passes 1--3 (−6.5, −8.6 and −11.1 points; post-confirmation,
simultaneous intervals excluding zero), so the pattern is pass-dependent
at fixed depth. Replacing the 79 controls by the seven same-domain names
lowers the gain to 20.8 points (9.1--32.5), because the lens also
recovers same-domain names about nine times more often than arbitrary
ones. An exhaustive tokenized-input audit finds no intended-label token
in any prompt and six control-token mentions whose removal moves the
estimate by at most 0.3 points. Changing the derivative target or the
position aggregation substantially reduces the late advantage. FP64
reference scoring and a cross-hardware state replay change the primary
estimate by at most 0.42 points. On the discovery set, retained readouts
of application-era banks that target each pass's own exit recover more
early-pass content than both the final-target lens and the raw lens; the
corresponding test on the confirmation set is blocked because those
banks survive only in a remote repository. The result supports a
localized, estimator- and target-specific improvement in concept
recovery, not general J-Lens superiority, improved answer accuracy, or
causal use of the recovered content.

\subsection{1. The question}\label{the-question}

A readout can miss useful content even when that content is present in a
model's hidden state. The raw logit lens applies the model's output head
directly to an intermediate activation. J-Lens instead first applies a
fixed, Jacobian-derived map that averages how activations affect
downstream states across calibration contexts, rather than fitting a
predictor to one task's answers. Its intended target is verbalizable
content, not accurate next-token prediction in every context {[}1{]}.

Ouro gives this comparison two depth axes: physical layer and recurrent
pass. The same layer stack is reused, so a readout that works at one
pass need not work at another. The experiment asks where each method
recovers a task-annotated intermediate concept more reliably. It does
not assume that a recovered word was part of the model's actual
reasoning {[}2{]}.

My initial expectation was that J-Lens would provide a broadly better
view of intermediate computation. That expectation did not survive the
early-pass results. The positive finding came later: a specific region
of the final pass where J-Lens outperformed raw readout. This paper
separates the exploration that found that region from the new-question
test that checked it, and reports the post-confirmation reviewer
analyses separately from both.

The distinction matters beyond this comparison. Hidden-state readouts
could complement observation of visible reasoning, but only after
establishing what their scores measure. No monitoring system, deception
detector, or chain-of-thought baseline is evaluated here.

\subsection{2. Readouts and measurement}\label{readouts-and-measurement}

\subsubsection{2.1 Model and fixed
estimators}\label{model-and-fixed-estimators}

The model is \textbf{ByteDance/Ouro-2.6B}, revision
\texttt{1ed04250da1a9936042725d302e81c8fa2ab5abd} (not the Thinking or
RLTT checkpoint), evaluated frozen in BF16 with four passes through 48
blocks, residual width 2,048 and a vocabulary of 49,152 tokens. I use
one-based physical layers \(\ell\) and pass numbers \(r\); the stored
zero-based virtual index is \(v = 48(r-1) + (\ell-1)\). Ouro
re-normalizes the residual stream at the end of every pass before
carrying it into the next pass and before the shared output head, so a
recorded state at the last block of a pass is the pre-normalization
block output.

For a captured residual state \(h_{i,r,\ell}\) at the readout position
of item \(i\), let \(N\) be the model's final RMS normalization and
\(W_U\) its unembedding. The two vocabulary readouts are

\[
z^{\mathrm{raw}}_{i,r,\ell}=W_U N(h_{i,r,\ell}), \qquad
z^{J}_{i,r,\ell}=W_U N(B_{r,\ell}h_{i,r,\ell}),
\]

where \(B_{r,\ell}\) is the frozen J-Lens matrix for that virtual layer.
Ranking the logits gives the token readout; no task classifier is
trained on the confirmation questions. The main banks target the final
residual state of pass 4 (virtual layer 191, before the final
normalization). The original J-Lens study's default Sonnet configuration
targets the penultimate layer, so this is an adaptation, not an exact
reproduction of that configuration {[}1{]}.

\textbf{Fitting recipe (from the executed records).} Each bank averages,
over 100 calibration paragraphs, the Jacobian estimator of the released
implementation: for every source layer \(v\) and output dimension, a
one-hot cotangent is placed at every valid target position of the target
layer at once and back-propagated through the BF16 model with ordinary
autograd, giving at source position \(p\) the sum over later positions
\(t' \ge p\) of \(\partial h_{191}[t']/\partial h_v[p]\); the mean over
valid source positions (positions 16 to len−2 of a 128-token input, BOS
included) is the paragraph's matrix. Matrices are accumulated as an
equal-paragraph FP32 sum, divided by 100 once and stored in FP16; at
readout they are cast back to FP32. Calibration paragraphs for
fit01--fit05 are independent random samples of 100 WikiText-103 training
paragraphs (at least 600 characters; seeds 2026090701--05); the realized
paragraph sets are pairwise disjoint and article overlap is 0--3 pairs
per fit pair. \textbf{fit01} was preselected for confirmation before any
new-item outcome; fit02 was retained as a secondary comparison. Fitting
ran on an RTX 5090 with an 8-direction CUDA-graph engine whose output
was verified bit for bit against the released fitter. Appendix C and the
accompanying methods ledger give every setting with its source.

Two further comparisons examine estimator specification. The
\textbf{target control} changes the derivative target from the final
block (virtual 191) to the penultimate block (190) while retaining final
normalization and unembedding. The \textbf{position controls} place the
cotangent at one uniformly drawn valid position \(q\) per paragraph and
keep all 2,048 derivative directions: \emph{sampled-sum} sums the
gradient over all valid source positions (its expectation over \(q\)
equals the dense estimator), whereas \emph{diagonal} keeps only the
source position \(p=q\). These are distinct estimator constructions, all
fitted on fit01's calibration paragraphs.

\subsubsection{2.2 Prompt, readout position and what counts as
recovery}\label{prompt-readout-position-and-what-counts-as-recovery}

Each question is one unfinished ``Fact:'' statement ending in a space;
the intermediate concept is implied by the clue and never named, and the
answer is the continuation. The model input is BOS plus the prompt
tokens up to the boundary where the prompt's tokenization stops agreeing
with the tokenization of prompt-plus-answer (in 131 of 160 items the
trailing-space token merges into the answer's first token and is
dropped); the readout position is the last input token, so every input
token is causally visible to the scored state. There is no instruction,
chat template, demonstration, prefilled response or padding; inputs are
13--35 tokens. Example (confirmation-001): ``Fact: The planet whose
sidereal rotation takes longer than its orbital period has an atmosphere
dominated by'' → readout at '' by'' (token 20), intermediate
\textbf{Venus} (scored token 16293, '' Venus''), answer ``carbon
dioxide'', never shown.

A concept is recovered when at least one accepted single-token form (the
name with and without a leading space, capitalized or lower-cased, 1--4
tokens per name) ranks among the first ten tokens of the \textbf{full
vocabulary}. Each item's controls are the other 79 intermediate names,
equally weighted. Let \(H_{i,r,\ell}^{m}(a)\) be the hit indicator for
name \(a\) under method \(m\), and \(a_i\) the intended name. With
control set \(C_i\), the fixed-layer excess score is

\[
e_{i,r,\ell}^{m}=H_{i,r,\ell}^{m}(a_i)
-\frac{1}{79}\sum_{c\in C_i}H_{i,r,\ell}^{m}(c).
\]

The primary endpoint is the equal-weight item mean of the paired
difference, averaged across the twelve fixed layers of the fourth pass:

\[
\widehat{\Delta}
=\frac{1}{160}\sum_{i=1}^{160}\frac{1}{12}
\sum_{\ell=26}^{37}\left(e_{i,4,\ell}^{J}-e_{i,4,\ell}^{\mathrm{raw}}\right).
\]

This averages fixed-layer measurements; it does \textbf{not} ask whether
a concept appears anywhere in the band. The historical \emph{any-layer}
score separately gives each intended name and each control 48
opportunities per pass, then averages controls, labels within items, and
items; it is reported separately.

All effects below are in percentage points. Recovery rates are averages
over item--layer cells, not the fraction of questions the model answers
correctly. Subtracting control recovery makes the comparison more
informative, but does not turn it into a causal test or a calibrated
measure of semantic correctness.

\subsection{3. From discovery to a frozen new-question
test}\label{from-discovery-to-a-frozen-new-question-test}

\subsubsection{3.1 What the discovery work
established}\label{what-the-discovery-work-established}

The discovery evaluation contained 90 eligible multihop items (100
labels) and 51 arithmetic items from the released lens-eval sets.
Fixed-layer analysis showed that the early multihop loss was not merely
an artifact of allowing the raw lens to recover a concept at a
favourable late layer: the disadvantage remained in the middle of the
stack. Arithmetic was mixed.

Pass 4 differed. Across five independent calibration fits, the mean
excess advantage in layers 26--37 was \textbf{18.72 points} (between-fit
SD \textbf{0.39}), positive in every fit, and every fit peaked at
physical layer 32 (Figure 1). The final third (layers 33--48) averaged
6.13 points (SD 0.21). The gain predominantly reflected increased
intended-concept recovery (18.91 own-name minus 0.19 control points).

That was evidence about calibration variability on the \textbf{same
questions}, not an independent replication across evaluation
populations. The band's crossed fit/item pointwise interval was
10.96--26.99 points, but its simultaneous interval over the 56
prespecified metrics (−0.86 to 38.29) and its shared-concept interval
(−7.49 to 44.92) contained zero, and no positive individual layer
cleared the simultaneous shared-concept band. The band was therefore
treated as a discovery to test, not a settled advantage.

\begin{figure}
\centering
\pandocbounded{\includegraphics[keepaspectratio,alt={Discovery population: fit01--fit05 minus raw, fixed-layer excess hit@10 for the 90 eligible multihop items, all four passes. Thin lines are the five fits, the bold line their mean, shading the crossed fit/item pointwise 95\% interval (not simultaneous). Yellow: layers 26--37, selected here and later frozen; grey: final third. The open point at layer 48 is the identity endpoint. Source: refit analysis arrays.}]{figures/fig1_discovery_five_fit_curves.png}}
\caption{Discovery population: fit01--fit05 minus raw, fixed-layer
excess hit@10 for the 90 eligible multihop items, all four passes. Thin
lines are the five fits, the bold line their mean, shading the crossed
fit/item pointwise 95\% interval (not simultaneous). Yellow: layers
26--37, selected here and later frozen; grey: final third. The open
point at layer 48 is the identity endpoint. Source: refit analysis
arrays.}
\end{figure}

\subsubsection{3.2 Confirmation population and
inference}\label{confirmation-population-and-inference}

The confirmation set contains \textbf{160 two-hop questions, 80
intermediate concepts, two questions per concept}, in ten domains of
eight concepts (arts and music, astronomy, biology, chemistry,
countries, environment and geology, historical people, landforms, SI
units, technology). All 160 were eligible; none of the 80 intermediates
or their aliases appears in the discovery set. Construction and
annotation were blind to method outcomes; factual hops, novelty and
control hygiene were reviewed independently before the local freeze.

The relation mixture differs from discovery: 96 questions request
relation types absent from the discovery inventory, 45 use familiar
requested relations, 19 recombine familiar relations. This is a transfer
test, not a difficulty-matched replication. Three component facts
overlap known calibration text (items 076, 130, 151). No claim is made
that the facts were absent from pretraining.

Fit01, the pass-4 band, scoring rules, population and analysis were
fixed before any new-item outcome; the stopping rule was one evaluation
of that population. ``Prospective'' refers to that recorded local freeze
(9 September 2026, 22:40:58 UTC); a public preregistration does not
exist. Two runtime corrections frozen before scoring and the later
numerical checks are described in Appendix B.

The primary interval uses \textbf{20,000 bootstrap resamples of whole
dependency groups} (seed 2026090901), with the 2.5th and 97.5th
percentiles as limits; resampled groups contribute all their items and
the ratio of sums preserves equal item weights. Twenty prespecified
secondary contrasts form a separate simultaneous 95\% family (bootstrap
max-\(t\)). Component intervals are descriptive.

\textbf{Dependency groups.} Items are joined into a group when they
share an intermediate (or its aliases), an underlying fact, a concrete
entity-substitution template (the same typed clue → entity →
requested-slot construction, such as element clue → element → atomic
number), or a normalized final answer or alias code. The 28 connected
components have sizes 42, 24, 16, 10, 10, 6, four of 4 and eighteen of
2; the weight-based effective group count is
\(160^2/\sum_g n_g^2 = 8.63\). The largest group joins the 16 chemistry
items (element→atomic-number and element→symbol templates) and the 16
SI-unit items (unit→symbol and unit→quantity templates) through
tungsten's and Watt's shared answer ``W'', with three small clusters
(Alps/Andes/Borneo; nucleus; ozone) attached by the small-integer
answers ``2'' and ``3''. Leaving out any single group keeps the estimate
between 19.73 and 24.87 points (low end without the 16 country items,
high end without the 42-item group). The effective count is a weight
summary, not a literal number of independent observations, and the
grouping does not certify the absence of other dependence.

\subsection{4. Results on new questions}\label{results-on-new-questions}

\subsubsection{4.1 The fixed final-pass advantage
transfers}\label{the-fixed-final-pass-advantage-transfers}

In pass 4, layers 26--37, fit01 exceeds the raw lens by \textbf{23.19
points} in excess hit@10, with a 95\% dependency-group interval of
\textbf{16.29--32.90 points} (Figure 2, Table 1).

{\def\LTcaptype{none} % do not increment counter
\begin{longtable}[]{@{}
  >{\raggedright\arraybackslash}p{(\linewidth - 4\tabcolsep) * \real{0.3333}}
  >{\raggedleft\arraybackslash}p{(\linewidth - 4\tabcolsep) * \real{0.3333}}
  >{\raggedright\arraybackslash}p{(\linewidth - 4\tabcolsep) * \real{0.3333}}@{}}
\toprule\noalign{}
\begin{minipage}[b]{\linewidth}\raggedright
Quantity
\end{minipage} & \begin{minipage}[b]{\linewidth}\raggedleft
Estimate
\end{minipage} & \begin{minipage}[b]{\linewidth}\raggedright
95\% group-bootstrap interval
\end{minipage} \\
\midrule\noalign{}
\endhead
\bottomrule\noalign{}
\endlastfoot
J-Lens intended recovery & 37.92\% & 30.44--46.56\% \\
Raw intended recovery & 14.43\% & 10.45--17.87\% \\
J-Lens control recovery & 0.34\% & 0.18--0.51\% \\
Raw control recovery & 0.04\% & 0.02--0.08\% \\
\textbf{Paired excess difference} & \textbf{+23.19 points} &
\textbf{+16.29 to +32.90 points} \\
\end{longtable}
}

\emph{Table 1. Confirmation results for the frozen pass-4 band.
Component intervals are descriptive; the final row is the primary
inference. Rounded entries need not subtract exactly.}

The intended-recovery difference is \textbf{23.49 points} (16.62--33.16,
descriptive). Control recovery is slightly higher under J-Lens (+0.30
points), reducing rather than creating its advantage: the excess comes
from recovering more intended concepts. Resampling individual questions
instead gives 17.97--28.71 points.

Fit02 gives a secondary advantage of \textbf{22.56 points} (simultaneous
interval 9.52--35.59); fit01 minus fit02 is 0.63 points (−0.73 to
+1.99). The two fixed estimators perform similarly here; this is not a
formal equivalence result. Discovery and confirmation results are not
pooled.

\begin{figure}
\centering
\pandocbounded{\includegraphics[keepaspectratio,alt={Confirmation population: fixed-layer curves for all four passes. Top: fit01 − raw (solid) and fit02 − raw (dashed) paired excess difference per layer, with the descriptive group-bootstrap pointwise 95\% band for fit01. Bottom: intended recovery for fit01 and raw, with control recovery dotted. Yellow: the frozen primary band. Only the band mean carries the prospective inference. Source: accepted payload readouts.}]{figures/fig2_confirmation_curves.png}}
\caption{Confirmation population: fixed-layer curves for all four
passes. Top: fit01 − raw (solid) and fit02 − raw (dashed) paired excess
difference per layer, with the descriptive group-bootstrap pointwise
95\% band for fit01. Bottom: intended recovery for fit01 and raw, with
control recovery dotted. Yellow: the frozen primary band. Only the band
mean carries the prospective inference. Source: accepted payload
readouts.}
\end{figure}

\subsubsection{4.2 The early-pass deficit transfers, including in the
same physical
band}\label{the-early-pass-deficit-transfers-including-in-the-same-physical-band}

All six prespecified early-pass comparisons remain negative on the new
population (Table 2).

{\def\LTcaptype{none} % do not increment counter
\begin{longtable}[]{@{}lll@{}}
\toprule\noalign{}
Pass & Mean over fixed layers 1--48 & Any-layer recovery \\
\midrule\noalign{}
\endhead
\bottomrule\noalign{}
\endlastfoot
1 & −8.83 {[}−13.64, −4.01{]} & −44.41 {[}−59.03, −29.78{]} \\
2 & −14.86 {[}−19.54, −10.18{]} & −57.37 {[}−72.01, −42.72{]} \\
3 & −15.59 {[}−22.40, −8.77{]} & −61.08 {[}−79.59, −42.58{]} \\
\end{longtable}
}

\emph{Table 2. Fit01 minus raw excess hit@10, in percentage points.
Brackets are simultaneous 95\% intervals from the 20-contrast secondary
family.}

These summaries alone do not show what happens in the specific band that
carries the pass-4 gain. As a \textbf{post-confirmation reviewer
analysis}, the same physical layers 26--37 were evaluated in every pass
with the frozen estimator, controls and groups (Table 3, Figure 3). The
band shows a deficit in each early pass whose simultaneous interval (a
four-contrast family, Appendix B) excludes zero; in pass 1 fit01
recovers no intended label in the band at all. The pattern is therefore
pass-dependent at fixed physical depth, not an artifact of averaging
over 48 layers.

{\def\LTcaptype{none} % do not increment counter
\begin{longtable}[]{@{}
  >{\raggedright\arraybackslash}p{(\linewidth - 10\tabcolsep) * \real{0.1667}}
  >{\raggedleft\arraybackslash}p{(\linewidth - 10\tabcolsep) * \real{0.1667}}
  >{\raggedleft\arraybackslash}p{(\linewidth - 10\tabcolsep) * \real{0.1667}}
  >{\raggedleft\arraybackslash}p{(\linewidth - 10\tabcolsep) * \real{0.1667}}
  >{\raggedright\arraybackslash}p{(\linewidth - 10\tabcolsep) * \real{0.1667}}
  >{\raggedright\arraybackslash}p{(\linewidth - 10\tabcolsep) * \real{0.1667}}@{}}
\toprule\noalign{}
\begin{minipage}[b]{\linewidth}\raggedright
Pass
\end{minipage} & \begin{minipage}[b]{\linewidth}\raggedleft
J-Lens intended
\end{minipage} & \begin{minipage}[b]{\linewidth}\raggedleft
Raw intended
\end{minipage} & \begin{minipage}[b]{\linewidth}\raggedleft
Paired excess difference
\end{minipage} & \begin{minipage}[b]{\linewidth}\raggedright
Interval
\end{minipage} & \begin{minipage}[b]{\linewidth}\raggedright
Status
\end{minipage} \\
\midrule\noalign{}
\endhead
\bottomrule\noalign{}
\endlastfoot
1 & 0.00\% & 6.46\% & −6.50 & {[}−11.91, −1.08{]} & post-confirmation,
simultaneous \\
2 & 4.22\% & 12.71\% & −8.64 & {[}−12.43, −4.85{]} & post-confirmation,
simultaneous \\
3 & 2.76\% & 13.91\% & −11.14 & {[}−16.76, −5.52{]} & post-confirmation,
simultaneous \\
4 & 37.92\% & 14.43\% & +23.19 & {[}+16.29, +32.90{]} & prospective
primary \\
\end{longtable}
}

\emph{Table 3. The same physical band, layers 26--37, in every pass.
Control recovery is below 0.35\% in every cell.}

\begin{figure}
\centering
\pandocbounded{\includegraphics[keepaspectratio,alt={Layers 26--37 in every pass on the confirmation population. Left: paired fit01 − raw excess difference; passes 1--3 carry the post-confirmation simultaneous 95\% interval, pass 4 the original prospective percentile interval. Right: intended recovery by method.}]{figures/fig3_same_band_by_pass.png}}
\caption{Layers 26--37 in every pass on the confirmation population.
Left: paired fit01 − raw excess difference; passes 1--3 carry the
post-confirmation simultaneous 95\% interval, pass 4 the original
prospective percentile interval. Right: intended recovery by method.}
\end{figure}

\subsubsection{4.3 The estimator specification
matters}\label{the-estimator-specification-matters}

The control comparisons transfer to the new questions (Table 4, Figure
4).

{\def\LTcaptype{none} % do not increment counter
\begin{longtable}[]{@{}
  >{\raggedright\arraybackslash}p{(\linewidth - 4\tabcolsep) * \real{0.3333}}
  >{\raggedleft\arraybackslash}p{(\linewidth - 4\tabcolsep) * \real{0.3333}}
  >{\raggedright\arraybackslash}p{(\linewidth - 4\tabcolsep) * \real{0.3333}}@{}}
\toprule\noalign{}
\begin{minipage}[b]{\linewidth}\raggedright
Contrast in pass 4, layers 26--37
\end{minipage} & \begin{minipage}[b]{\linewidth}\raggedleft
Difference
\end{minipage} & \begin{minipage}[b]{\linewidth}\raggedright
Simultaneous 95\% interval
\end{minipage} \\
\midrule\noalign{}
\endhead
\bottomrule\noalign{}
\endlastfoot
Main fit01 minus penultimate target & +14.38 & {[}+7.70, +21.07{]} \\
Sampled-sum minus diagonal & +15.77 & {[}+5.80, +25.75{]} \\
Penultimate target minus raw & +8.80 & {[}+0.76, +16.85{]} \\
Sampled-sum minus raw & +21.72 & {[}+9.16, +34.29{]} \\
Diagonal minus raw & +5.95 & {[}−0.34, +12.24{]} \\
\end{longtable}
}

\emph{Table 4. Estimator controls in percentage points; matched layer
support within each control family; same 20-contrast family as Table 2.}

Changing the derivative target by one block reduces the advantage
substantially; restricting the position reduction to the matching source
position does so as well. Neither comparison isolates how Ouro reasons:
it shows that the measured recovery depends on the lens's construction.
The diagonal-only estimate is smaller and its support is borderline: the
lower interval bound is within 0.4 points of zero on every numerical
path and excludes zero only under regenerated states. The direct
sampled-sum-minus-diagonal comparison is the clearer result.

\begin{figure}
\centering
\pandocbounded{\includegraphics[keepaspectratio,alt={Pass-4 band contrasts. Black: the prospective primary (percentile interval). Grey: contrasts of the original 20-member simultaneous family. Blue: the post-confirmation within-domain-control contrast (Section 4.4).}]{figures/fig4_band_contrasts.png}}
\caption{Pass-4 band contrasts. Black: the prospective primary
(percentile interval). Grey: contrasts of the original 20-member
simultaneous family. Blue: the post-confirmation within-domain-control
contrast (Section 4.4).}
\end{figure}

\subsubsection{4.4 Post-confirmation reviewer
checks}\label{post-confirmation-reviewer-checks}

Three further analyses were specified on the already-used confirmation
data after the confirmation; they are reported separately and do not
alter the primary endpoint.

\textbf{Within-domain controls.} Replacing each item's 79 controls by
the seven other names in its benchmark domain gives a paired excess
difference of \textbf{20.80 points} (simultaneous 95\% interval
9.13--32.46). J-Lens recovers same-domain names at 3.07\% against 0.34\%
for arbitrary controls (raw: 0.38\% against 0.04\%), so about a ninth of
the lens's gain is domain- or type-level rather than concept-specific.
An arithmetic bound derived from the published values (intended
difference minus 79/7 times the J-Lens all-control rate, 19.61 points)
is satisfied by the exact arrays for every item. Four of the ten domains
contain a single semantic type (elements, countries, SI units, people);
the others mix types (planets and moons; rivers, ranges, a desert and an
island; molecules, a cell type and an organelle), so this is a
within-domain control, not a same-type control. The gain is
heterogeneous across domains: from about +51 points for countries and
+40 for environment and geology to about zero for SI units, where
neither method recovers unit names in this band, and to −1 for astronomy
(+5 with the 79 controls), where the lens often recovers other planets
and moons. Domains are not independence units.

\textbf{Tokenized-input overlap.} For every item, every scored token of
the intended label and of all 79 controls was compared with the exact
model input. No intended-label token appears in any input. Six control
tokens do, in six items: ``Voyager'' in three astronomy prompts,
``Chile'' in an Andes prompt, ``Egypt'' in a Nile prompt, and the
lower-case form ``io'' as a token inside ``Fidelio''. Dropping those
controls for the affected items (both methods; 78 controls for six
items) gives 23.19 points (16.25--32.72); dropping the six items gives
22.90 points (15.92--33.15). Lexical word-level mentions without
scored-token overlap occur once more (``newton'' in a pascal item). Zero
token overlap does not establish absence of semantic clues or of the
facts from pretraining; the prompts are designed to imply the
intermediate.

\textbf{Grouping.} The rule, memberships and the composition of the
42-item group are given in Section 3.2 and the reviewer ledger; no
alternative grouping was constructed.

\subsection{5. Numerical verification}\label{numerical-verification}

The confirmation run used BF16 model execution on an RTX 5090. A
development check initially failed because applying the output head to
one residual row did not reproduce logits computed in a multi-row call.
The subsequent verification measured whether this affected the scored
comparisons rather than assuming a small logit discrepancy was harmless.
These checks were \textbf{post-confirmation}, not prospectively
registered.

On a local RTX 5070 Ti with the confirmation's torch build and precision
settings, the retained confirmation states, verified banks and frozen
scoring code reproduced the saved arrays bit for bit across \textbf{160
items, six arms and all scored columns}, including ranks, top-ten token
IDs, saved exit logits and sampled transported vectors. Calls with 160,
190, 191 and 192 rows matched the executed layout exactly; single-row
calls changed 1.5--1.6 million of about 1.5 billion logits per arm by at
most 0.125 (0.25 for raw) but no endpoint quantity.

{\def\LTcaptype{none} % do not increment counter
\begin{longtable}[]{@{}
  >{\raggedright\arraybackslash}p{(\linewidth - 6\tabcolsep) * \real{0.2500}}
  >{\raggedleft\arraybackslash}p{(\linewidth - 6\tabcolsep) * \real{0.2500}}
  >{\raggedright\arraybackslash}p{(\linewidth - 6\tabcolsep) * \real{0.2500}}
  >{\raggedleft\arraybackslash}p{(\linewidth - 6\tabcolsep) * \real{0.2500}}@{}}
\toprule\noalign{}
\begin{minipage}[b]{\linewidth}\raggedright
Computation
\end{minipage} & \begin{minipage}[b]{\linewidth}\raggedleft
Primary difference
\end{minipage} & \begin{minipage}[b]{\linewidth}\raggedright
95\% group interval
\end{minipage} & \begin{minipage}[b]{\linewidth}\raggedleft
Change from accepted
\end{minipage} \\
\midrule\noalign{}
\endhead
\bottomrule\noalign{}
\endlastfoot
Accepted BF16 path & +23.19 & {[}+16.29, +32.90{]} & --- \\
Alternative row packing & +23.19 & unchanged & 0.00 \\
FP64 reference computation & +22.93 & {[}+16.17, +32.52{]} & −0.26 \\
Locally regenerated states & +23.61 & {[}+16.52, +33.77{]} & +0.42 \\
\end{longtable}
}

\emph{Table 5. Numerical sensitivity of the primary endpoint. The FP64
row applies a reference readout computation to the stored tensors, not
an FP64 model. Alternative packing was measured on the local GPU, not in
a fresh RTX 5090 run.}

CPU and GPU FP64 computations gave identical ranks, hits and top-ten
sets (values differing by at most \(9\times10^{-14}\)). Resolving every
exactly tied logit adversarially bounds the accepted estimate between
\textbf{22.60 and 23.35 points}; moving every logit by up to 1/32 or 1/8
in the adversarial direction gives 22.01--24.41 and 18.30--27.26.
Regenerating states from the prompts on the local GPU was not
bit-identical to the RTX 5090 forward: 37 of 160 state records matched
exactly and 154 greedy continuations matched, yet replaying every arm on
those states moved the endpoint by only 0.42 points. All six early-pass
contrasts and the direct estimator-control contrasts keep their signs
and zero-exclusion on every path. Bit-level state regeneration on
another RTX 5090 was not tested.

A separately implemented checker (a separate coding agent working from a
specification, without producer code) reconstructed the saved analysis
to within \(9\times10^{-16}\) and detected all 13 planted errors in
isolated fixtures. ``Independent'' here means a separate implementation,
not replication by another laboratory.

\subsection{6. Local-exit targets: a conditional comparison, partly
blocked}\label{local-exit-targets-a-conditional-comparison-partly-blocked}

The main bank targets the final pass, so early-pass states are far from
its target: a pass-1 state at layer 30 is 162 blocks and three
pass-boundary normalizations away from virtual layer 191. One
explanation of the early deficit is therefore transport distance rather
than a property of early-pass representations. The intended test
evaluates the same stored early-pass confirmation states with the raw
lens, a J-Lens targeting the current pass's own exit, and the
final-target J-Lens, in layers 26--37 of passes 1--3.

\textbf{Inventory.} Local-exit banks exist only from the application era
(5 September 2026): banks targeting virtual 47, 95, 143 and 191, fitted
on the first 100 paragraphs of the historical calibration file (12
articles) with one shared recipe that differs only in the target. Their
merged files were purged; the final-target bank and four of its five
shards survive locally, whereas the twelve shards of the three
local-exit banks survive only in a private repository (hashes recorded,
no credential on this machine). The confirmation-population comparison
is therefore \textbf{blocked}, and no comparison was improvised.
Retrieval alone (11 GB, no fitting, no paid compute) would unblock it; a
fit01-matched refit would cost about eight GPU-hours and would be a new
estimator family.

\textbf{What the retained readouts show (descriptive, discovery
population).} The application-era banks were all applied to the 148
discovery items in one evaluation whose rank arrays survive. On the 90
eligible multihop items, in the fixed band 26--37, the lens targeting
each pass's own exit recovers the intended concept far more often than
both alternatives: intended recovery 15.4\%, 16.2\% and 25.2\% in passes
1--3, against 2.1\%, 4.3\% and 10.1\% for the raw lens and 0.3\%, 0.4\%
and 2.1\% for the final-target lens (Figure 5). The paired excess
differences, local minus final, are +14.6, +16.9 and +23.0 points
(shared-concept-cluster percentile intervals 9.1--20.9, 11.5--22.8 and
16.0--30.6), and local minus raw +13.0, +11.3 and +14.9 points
(8.1--19.0, 6.6--16.6, 7.9--22.4). Control recovery stays below 2.3\%
everywhere.

\begin{figure}
\centering
\pandocbounded{\includegraphics[keepaspectratio,alt={Application-era matched banks (same recipe, target differs) evaluated on the DISCOVERY population, fixed layers, passes 1--3: raw lens, final-target J-Lens (exit 3) and the pass's own-exit J-Lens. Right: band means with descriptive shared-concept-cluster 95\% intervals. Post hoc and descriptive; not the confirmation population.}]{figures/fig5_local_exit_discovery_descriptive.png}}
\caption{Application-era matched banks (same recipe, target differs)
evaluated on the DISCOVERY population, fixed layers, passes 1--3: raw
lens, final-target J-Lens (exit 3) and the pass's own-exit J-Lens.
Right: band means with descriptive shared-concept-cluster 95\%
intervals. Post hoc and descriptive; not the confirmation population.}
\end{figure}

\textbf{Interpretation.} Within the rules fixed before these numbers
were computed, this is the case in which the local target beats both
alternatives: a target-dependent early-pass advantage, consistent with a
transport explanation. It is not proof that averaging destroyed
information; it is on the population on which the band was selected; it
uses the older 12-article calibration and a different batch setting; the
intervals are descriptive; and it has not been tested on the
confirmation population. Two existing qualifications remain: shortening
the final target by one block (the penultimate control) reduced the late
gain, so transport distance alone is not established, and raw pass
counts are not comparable transport lengths across architectures. The
historical Huginn pilot (Appendix A) did not identify a cause or rule
out supervision-based explanations.

\subsection{7. Interpretation and
limits}\label{interpretation-and-limits}

The supported result is a \textbf{localized improvement in reading
annotated intermediate concepts from fixed Ouro states}: a specific
final-pass band, one frozen estimator, transferred to new questions and
concepts without refitting or band selection. The same band in the
earlier passes shows the opposite sign, so the method is not uniformly
better as computation proceeds, and the source of the measured gain is
higher intended-concept recovery, part of which is domain-level.

One possible explanation for the strong raw baseline is Ouro's training
objective, which applies weighted next-token losses at recurrent pass
boundaries through a shared output head {[}2{]}. This does not supervise
every physical block, and the present experiments do not isolate the
effect of that objective. The target-dependence seen on the discovery
population points to transport distance or target choice as a
contributing factor, but architecture, calibration, target choice and
executed depth remain entangled, and the Huginn pilot's evidence is
limited by its incomplete artifact chain.

Recovery is not use. In exploratory checks a lens recovered a labelled
intermediate on questions the model answered incorrectly, and native
exits can agree while their readouts disagree sharply (Appendix A).

The main population limit is the small number of unequal dependency
groups. The interval accounts for the declared grouping, not every
possible shared fact, relation or template; the effective count of 8.63
warns against reading 160 related questions as 160 demonstrations. The
relation mixture is not difficulty-matched, labels are restricted to
accepted single-token forms, the domain heterogeneity is large, and the
evidence concerns one checkpoint and two retained calibration fits. The
study does not establish performance on arbitrary concepts, other
models, or monitoring under adversarial pressure.

\subsubsection{Reproducibility and artifact
availability}\label{reproducibility-and-artifact-availability}

The confirmation banks, states, scores, frozen inputs and code permit
exact table reconstruction, rescoring from banks and states, and forward
replay (not bit-identical off the RTX 5090). Complete refitting of the
original estimators is not available: an earlier retrieval failure and a
later storage purge destroyed unique historical artifacts, including all
four Ouro FP32 fitting checkpoints, the historical discovery-evaluation
state cache, the truncated Huginn checkpoint and the application-era
local-exit banks (43 purged paths, 33 distinct contents lost; 19 paths
unresolved in a remote repository). Knowing the recipe does not restore
those checkpoints, and the Huginn pilot is retained only at the
saved-output level. On 11 September 2026 a copy-only archive of every
surviving file (16,200 files, 52.15 GB, every byte re-hashed on the
destination) was written to one external SSD and restoration-tested from
that SSD alone with the laptop's paths hidden and no network: the frozen
code matched its freeze, all four bank files loaded, the frozen analysis
reproduced the primary and all twenty secondary contrasts exactly, and
raw and fit01 rescoring from the archived states and weights reproduced
the saved ranks bit for bit. Until a second destination exists it is a
single-copy archive, not a redundant backup. The frozen records, code,
compact results, figures and this manuscript are published in the
repository https://github.com/VykosMolt/JLens-Ouro (research directory;
files of 50 MB or more, including the banks and state caches, are listed
there with hashes but not stored). No DOI exists at the time of writing.

\subsection{8. Conclusion}\label{conclusion}

J-Lens loses to raw readout in Ouro's first three passes, including in
the very layers where it wins in the fourth, and recovers more annotated
intermediate content in a fixed region of the final pass. That advantage
transferred to new questions and concepts and survives the measured
numerical variations and control redefinitions. Its size depends
strongly on the derivative target and the position aggregation, and
retained discovery-set readouts suggest that a lens targeting the pass's
own exit reads early passes far better than the final-target lens. The
finding is about where a particular readout works, not about general
superiority, successful reasoning or causal control.

\newpage

\subsection{Appendix A. Exploratory findings kept separate from
confirmation}\label{appendix-a.-exploratory-findings-kept-separate-from-confirmation}

\subsubsection{Exit disagreement}\label{exit-disagreement}

In the 7 September analysis of 148 prompts, native pass-3 and final
predictions agreed \textbf{85.1\%} of the time (item-bootstrap interval
79.1--90.5\%). Two J-Lens variants targeting the current and final exits
agreed only \textbf{1.45\%} (0.58--2.51\%) averaged over physical layers
41--47; layer 48 was excluded because some agreement there follows by
construction. One token earlier, native agreement rose to 96.6\% while
lens agreement remained 0.39\%. These are different comparisons, so the
lens gap cannot be read as the amount of model revision, and poor output
prediction does not by itself refute an intermediate-recovery method
{[}1{]}.

\subsubsection{Passing continuations and
arithmetic}\label{passing-continuations-and-arithmetic}

The early multihop deficit remained when the discovery items were
separated by the historical continuation-pass criterion. That criterion
was imperfect: two numerical false positives reduced the original 72
passes out of 148 to 70, while additionally accepting the numeric
equivalent of one word-form target gave 71. Four-token continuations and
lexical matching do not establish semantic correctness; the subgroup
observations are exploratory.

A separate supervised arithmetic probe scored 59.4\% against the raw
lens's 76.4\% in pass 1, but exceeded it in pass 3 (35.6\% versus
17.9\%). The classifiers trained on only 36 independent operand pairs; a
bounded training-size comparison (+3.9 points, interval −2.9 to +10.9)
was inconclusive. This used a different task family and a candidate-set
metric, with incomplete provenance for the older J-Lens comparator; it
is not pooled with the concept-recovery results.

\subsubsection{Huginn pilot}\label{huginn-pilot}

The pilot used one Huginn estimator calibrated on the same 100 texts as
the preselected Ouro fit, eight recurrent passes and two evaluation
initializations. Joint tokenizer eligibility retained 88 multihop items
(98 labels) and 51 arithmetic items. J-Lens-minus-raw estimates were
negative under all six prespecified summaries in both initializations;
the cross-model differences of differences had mixed signs and
simultaneous intervals containing zero. The prediction of a larger
relative J-Lens advantage in Huginn was not supported. This does not
isolate supervision: Huginn's coda is trained after sampled recurrence
depths {[}3{]}, so it is not an unsupervised-intermediate control, and
its raw-readout weakness is block-specific {[}4{]}. The fitted Huginn
bank was never retrieved and its partial checkpoint was later lost; the
saved statistics are historical evidence, not a verified cross-model
replication.

\newpage

\subsection{Appendix B. Verification and analysis
provenance}\label{appendix-b.-verification-and-analysis-provenance}

The original confirmation freeze is unchanged. Two separately recorded
corrections preceded new-item scoring: \texttt{precision\_v1} (10
September, 10:44:52 UTC) removed redundant TF32 setters, and
\texttt{native\_check\_v1} (22:56:18 UTC) changed the development
comparison from a single residual row to a whole-sequence readout and
resized run budgets. The checker's accepted run-specification pin was
updated after outcomes were known. The post-confirmation verification
reviewed these changes and found the scientific input, estimator,
endpoint and analysis unchanged.

The numerical results use the accepted BF16 endpoint as primary; FP64,
alternate packing, tie handling and regenerated states are sensitivity
checks. For diagonal minus raw, the accepted estimate is +5.95 points
with simultaneous interval −0.34 to +12.24; FP64 gives −0.40 as lower
bound and regenerated states +0.04, which is why its support is
described as borderline.

\textbf{Post-confirmation reviewer family (Section 4.2 and 4.4).}
Exactly four new inferential contrasts were fixed in a dated plan before
computation: the within-domain paired excess difference in the pass-4
band and the same-band paired differences in passes 1--3. They share one
centered bootstrap max-\(t\) 95\% family (20,000 whole-group draws, seed
2026091101, quantile 2.639). Every other new quantity is descriptive.
The overlap sensitivities and the local-exit discovery-population
summary are descriptive.

\textbf{Records.} Confirmation run
\texttt{ouro\_confirmation\_20260911\_fixed160\_native1}; frozen plan,
benchmark and freeze receipt under
\texttt{research/confirmation\_2026-09-09/}; verification under
\texttt{research/verification\_2026-09-11/}
(\texttt{FINAL\_VERIFICATION.md},
\texttt{REPORT\_confirmation\_verified.md},
\texttt{RUN\_SPECIFICATION.json}); this round's reviewer checks, methods
ledger, numerical source ledger and archive receipt under
\texttt{research/closeout\_2026-09-11/}. These are project provenance
identifiers, not public links.

\subsection{Appendix C. Method details in
brief}\label{appendix-c.-method-details-in-brief}

\begin{itemize}
\tightlist
\item
  \textbf{Tokenization and input.} BOS + prompt tokens
  (\texttt{truncation=True,\ max\_length=511}, never triggered); readout
  context = common prefix of encode(prompt) and encode(prompt+answer);
  position −1; batch size 1; SDPA attention; \texttt{use\_cache=False};
  BF16 weights.
\item
  \textbf{Hooks and normalization.} Forward hook on each shared physical
  block filtered by the recurrent step; recorded tensor = post-residual
  block output, before the pass-end RMS normalization; states stored as
  FP32 copies of BF16 values.
\item
  \textbf{Readout.} \(J h\) in FP32 (raw: \(h\)), cast to BF16,
  RMS-normalize (FP32 internally, \(\epsilon=10^{-6}\)), BF16 head, one
  call per item over all scored columns; descending argsort over 49,152
  tokens; rank = min over accepted forms; hit = rank \textless{} 10.
\item
  \textbf{Calibration and fitting.} 100 WikiText-103 paragraphs per fit
  (fit01--05 seeds 2026090701--05; paragraph-disjoint; 100/99/99/100/100
  distinct articles), inputs of 128 tokens including BOS (six shorter),
  valid positions 16 to len−2, cotangent summed over valid target
  positions and gradient averaged over valid source positions, FP32 sum
  over paragraphs divided by 100, FP16 storage, 8 output directions per
  backward pass with CUDA-graph replay (bit-parity verified), RTX 5090.
\item
  \textbf{Controls.} Penultimate: target virtual 190, sources 0--189.
  Positions: one uniform \(q\) per paragraph (seed 2026090801),
  sampled-sum = \(\sum_p \partial h_{191}[q]/\partial h_v[p]\) (no
  divisor), diagonal = \(\partial h_{191}[q]/\partial h_v[q]\), from one
  backward stream. Matched supports for contrasts: band 26--37 for all
  arms; final-third intersections 176--189 (target control) and 176--190
  (position controls).
\item
  \textbf{Banks.} fit01 \texttt{90f01f6a…}, fit02 \texttt{101f31db…}
  (byte-exact reconstruction from its intact checkpoint), penultimate
  \texttt{dc6354df…}, positions \texttt{b8e8b7d2…} (both position arms);
  four files, five bank arms, six evaluated arms.
\item
  \textbf{Application-era family (Section 6).} Same estimator with the
  first 100 historical paragraphs (12 articles), 32 directions per
  backward pass, B300 pod; targets 47, 95, 143, 191.
\end{itemize}

\subsection{Appendix D. Contribution and AI-assistance statement
(draft)}\label{appendix-d.-contribution-and-ai-assistance-statement-draft}

AI coding agents implemented the fitting, evaluation, controller and
analysis code, executed the runs under the author's authorization, and
performed separately instantiated review and numerical-reconstruction
passes; the project materials name Claude Code, Claude Opus and Codex
for the application era, and round logs after 7 September do not record
agent versions. The author set the research questions, approved the
frozen design and budgets, authorized each paid run and the storage
purge, and is responsible for all claims. Both manuscript drafts were
prepared by AI assistants from the frozen reports and records; the
extent of the author's verification of each number is to be stated by
the author. No venue policy has been assumed. {[}Items requiring the
author's confirmation are listed in the accompanying fact sheet.{]}

\subsection{References}\label{references}

\textbf{{[}1{]}} Gurnee, W., Sofroniew, N., Pearce, A., et
al.~\emph{Verbalizable Representations Form a Global Workspace in
Language Models.} Transformer Circuits, Anthropic, 6 July 2026. Methods;
Quantitative Methodological Comparisons; Methodological Details and
Ablations. https://transformer-circuits.pub/2026/workspace/

\textbf{{[}2{]}} Zhu, R.-J., Wang, Z., Hua, K., et al.~\emph{Scaling
Latent Reasoning via Looped Language Models.} arXiv:2510.25741, 2025;
Sections 3.1--3.3 for the architecture and training objective.
https://arxiv.org/abs/2510.25741

\textbf{{[}3{]}} Geiping, J., et al.~\emph{Scaling up Test-Time Compute
with Latent Reasoning: A Recurrent Depth Approach.} arXiv:2502.05171,
2025; Sections 3.1--3.3 and 6.1 for the Huginn architecture, training
and early exits. https://arxiv.org/abs/2502.05171

\textbf{{[}4{]}} Lu, W., et al.~\emph{Latent Chain-of-Thought? Decoding
the Depth-Recurrent Transformer.} arXiv:2507.02199, 2025; Section 3.2
for block-specific logit-lens readability in Huginn.
https://arxiv.org/abs/2507.02199

\end{document}
